% 数学ノート集・共通テンプレート v2（積分論2026年度第1部の版面を参考）
% A5判。章内で定義・定理・例・注意等の番号を共通化。
% コンパイル： platex -interaction=nonstopmode math_note_template_v2.tex (2回)
%              dvipdfmx math_note_template_v2.dvi
\documentclass[a5paper,10pt,onecolumn]{jsarticle}
\usepackage[dvipdfmx,top=22mm,bottom=20mm,inner=19mm,outer=18mm,headheight=13pt,headsep=6mm]{geometry}
\usepackage{amsmath,amssymb,amsthm,mathtools}
\usepackage{enumitem,fancyhdr,titlesec}
\usepackage[dvipdfmx,bookmarks=true,colorlinks=false]{hyperref}
\usepackage{pxjahyper}

% ----- 差し替える文書情報 -----
\newcommand{\NoteTitle}{平方剰余の相互法則}
\newcommand{\NoteSeries}{数学ノート集}

% ----- 文字・組版 -----
% 参考PDFの質素な明朝＋Computer Modernの印象を引き継ぎながら，
% A5判で窮屈にならないよう本文の行間を少し確保する．
\linespread{1.13}
\setlength{\parindent}{1zw}
\setlength{\parskip}{0pt}
\setlength{\abovedisplayskip}{6pt plus 2pt minus 2pt}
\setlength{\belowdisplayskip}{6pt plus 2pt minus 2pt}
\setlength{\abovedisplayshortskip}{3pt plus 2pt minus 1pt}
\setlength{\belowdisplayshortskip}{4pt plus 2pt minus 1pt}
\setlist[enumerate]{topsep=3pt,itemsep=1pt,parsep=0pt,leftmargin=2.4zw}
\titleformat{\section}{\normalfont\normalsize\bfseries}{\thesection}{0.8zw}{}
\titlespacing*{\section}{0pt}{2.0ex plus .4ex}{.9ex}
\titleformat{\subsection}{\normalfont\normalsize\bfseries}{\thesubsection}{0.7zw}{}
\titlespacing*{\subsection}{0pt}{1.5ex plus .3ex}{.6ex}
\raggedbottom

% ----- 柱とページ番号：ページ上部、罫線なし -----
\fancypagestyle{mathnote}{%
  \fancyhf{}%
  \fancyhead[L]{\footnotesize \NoteSeries\ （\NoteTitle）}%
  \fancyhead[R]{\footnotesize \thepage}%
  \renewcommand{\headrulewidth}{0pt}%
  \renewcommand{\footrulewidth}{0pt}%
}
\pagestyle{mathnote}

% ----- 節内通し番号の定義・定理群 -----
\newtheoremstyle{mathnotetheorem}
 {4.5pt}{4.5pt}{\normalfont}{}%
 {\bfseries}{}{0.65zw}%
 {\thmname{#1}\thmnumber{ #2}\thmnote{（#3）}}
\theoremstyle{mathnotetheorem}
\newtheorem{theorem}{定理}[section]
\newtheorem{definition}[theorem]{定義}
\newtheorem{lemma}[theorem]{補題}
\newtheorem{proposition}[theorem]{命題}
\newtheorem{corollary}[theorem]{系}
\newtheorem{example}[theorem]{例}
\newtheorem{remark}[theorem]{注意}
\renewcommand{\proofname}{証明}
\renewcommand{\qedsymbol}{$\square$}
\makeatletter
\renewenvironment{proof}[1][\proofname]{%
  \par\pushQED{\qed}\normalfont\topsep4\p@\@plus2\p@\relax
  \trivlist\item[\hskip\labelsep\bfseries #1]\ignorespaces
}{\popQED\endtrivlist\@endpefalse}
\makeatother

% ----- 使うノートで適宜拡張する記号 -----
\newcommand{\R}{\mathbb{R}}
\newcommand{\norm}[1]{\lVert #1\rVert}
\newcommand{\abs}[1]{\lvert #1\rvert}

\begin{document}
\section*{読了後にできること}
本稿を読むための前提は，Legendre記号，Eulerの規準，平方剰余相互法則．
本稿の具体的な到達課題は，奇素数に対する平方剰余性を相互法則と補充法則から判定する．

\section{平方剰余とLegendre記号}

以下，$p$ は奇素数とし，合同式の法は特に断らない限り $p$ とする．

\begin{definition}
$p\nmid a$ とする．合同式 $x^2\equiv a\pmod p$ が
整数解を持つとき $a$ を\emph{平方剰余}，
持たないとき\emph{平方非剰余}という．
Legendre記号を
\[
 \left(\frac ap\right)=
 \begin{cases}
 0&(p\mid a),\\
 1&(p\nmid a,\ a\text{ は平方剰余}),\\
 -1&(p\nmid a,\ a\text{ は平方非剰余})
 \end{cases}
\]
と定義する．
\end{definition}

\begin{lemma}\label{lem:squares}
$\mathbb F_p^\times$ に平方剰余はちょうど $(p-1)/2$ 個ある．
さらに平方剰余全体は
$\mathbb F_p^\times$ の指数2の部分群をなす．
\end{lemma}
\begin{proof}
$x^2=y^2$ であることと
$(x-y)(x+y)=0$ であることは同値である．
$\mathbb F_p$ は体で $p$ は奇数なので，
$x\ne0$ に対して $x$ と $-x$ は相異なる二つの元であり，
平方写像 $x\mapsto x^2$ の各像はちょうど二つの原像を持つ．
従って像の位数は $(p-1)/2$ である．
平方写像は乗法群の準同型なので，像は部分群であり，
その指数は2である．
\end{proof}

\begin{theorem}[Eulerの規準]\label{thm:euler}
任意の整数 $a$ に対し
\[
 \left(\frac ap\right)\equiv a^{(p-1)/2}\pmod p
\]
が成り立つ．
また $p\nmid ab$ なら
\[
 \left(\frac{ab}{p}\right)
 =\left(\frac ap\right)\left(\frac bp\right).
\]
\end{theorem}
\begin{proof}
$p\mid a$ の場合は両辺とも零である．
$p\nmid a$ とし，$m=(p-1)/2$ と置く．
Fermatの小定理により $a^{2m}=1$ なので，
体 $\mathbb F_p$ において $a^m\in\{1,-1\}$ である．
$a$ が平方 $x^2$ なら
$a^m=x^{p-1}=1$ である．
一方，多項式 $X^m-1$ の根は高々 $m$ 個であり，
補題\ref{lem:squares}によって平方剰余の $m$ 個が
すべてその根である．
従って平方非剰余なら $a^m=-1$ となる．
これでEulerの規準を得る．

$p\nmid ab$ のとき
$(ab)^m=a^mb^m$ なので，Eulerの規準により
Legendre記号の両辺は法 $p$ で合同である．
両辺はそれぞれ $\pm1$ であり，
奇素数 $p$ に対して $1\not\equiv-1\pmod p$ だから
整数として等しい．
\end{proof}

\section{Gaussの補題}

\begin{lemma}[Gaussの補題]\label{lem:gauss}
$p$ を奇素数，$m=(p-1)/2$ とし，$p\nmid a$ とする．
$aj$（$1\le j\le m$）の法 $p$ の最小正剰余のうち，
$p/2$ より大きいものの個数を $n(a)$ と置く．
このとき
\[
 \left(\frac ap\right)=(-1)^{n(a)}
\]
が成り立つ．
\end{lemma}
\begin{proof}
$aj$ の法 $p$ の最小絶対剰余を
$r_j$（$-m\le r_j\le m$, $r_j\ne0$）とする．
その絶対値 $|r_j|$ は $1,\dots,m$ の置換である．
実際，$|r_i|=|r_j|$ なら
$ai\equiv\pm aj\pmod p$ であり，
$p\nmid a$ より $i\equiv\pm j\pmod p$ となる．
$1\le i,j\le m$ だから $i=j$ である．

$aj\equiv r_j\pmod p$ を $j=1,\dots,m$ にわたり掛けると
\[
 a^m m!\equiv\prod_{j=1}^m r_j
 =(-1)^{n(a)}m!\pmod p.
\]
$p\nmid m!$ だから
$a^m\equiv(-1)^{n(a)}\pmod p$ を得る．
Eulerの規準と比較して結論が従う．
\end{proof}

\begin{lemma}[Eisensteinの補題]\label{lem:eisenstein}
$p$ を奇素数，$a$ を奇整数で $p\nmid a$ とする．
このとき
\[
 \left(\frac ap\right)
 =(-1)^{\sum_{j=1}^{(p-1)/2}\lfloor aj/p\rfloor}.
\]
\end{lemma}
\begin{proof}
まず $a>0$ とする．$m=(p-1)/2$ と置き，
$aj=pq_j+s_j$，
$q_j=\lfloor aj/p\rfloor$，
$1\le s_j\le p-1$ と書く．
$s_j>p/2$ となる $j$ の個数を $n(a)$ とし，
最小絶対剰余を $r_j$ と書く．
$1\le k_j=|r_j|\le m$ は $1,\dots,m$ の置換であり，
$s_j=k_j$ または $s_j=p-k_j$ だから
\[
 \sum_{j=1}^m s_j
 =\sum_{j=1}^mk_j
  +pn(a)-2\sum_{s_j>p/2}k_j.
\]
一方 $\sum_js_j=a\sum_jj-p\sum_jq_j$ なので
\[
 (a-1)\sum_{j=1}^m j
 =p\left(n(a)+\sum_{j=1}^m q_j\right)
  -2\sum_{s_j>p/2}k_j.
\]
$a-1$ は偶数で $p$ は奇数だから，
両辺を法2で比較すると
$n(a)\equiv\sum_jq_j\pmod2$ を得る．
Gaussの補題から結論が従う．

$a<0$ の場合には $a$ を正の奇数 $a+2kp$ に置き換える．
Legendre記号は変わらず，
\[
 \sum_{j=1}^m\left\lfloor\frac{(a+2kp)j}{p}\right\rfloor
 =\sum_{j=1}^m\left\lfloor\frac{aj}{p}\right\rfloor
  +2k\sum_{j=1}^m j
\]
なので指数の偶奇も変わらない．
従って一般の奇整数の場合にも成立する．
\end{proof}

\section{二つの補充法則}

\begin{theorem}[第一補充法則]\label{thm:first}
任意の奇素数 $p$ について
\[
 \left(\frac{-1}{p}\right)=(-1)^{(p-1)/2}
 =
 \begin{cases}
 1&(p\equiv1\pmod4),\\
 -1&(p\equiv3\pmod4).
 \end{cases}
\]
\end{theorem}
\begin{proof}
Eulerの規準に $a=-1$ を代入すれば
$\left(\frac{-1}{p}\right)\equiv(-1)^{(p-1)/2}\pmod p$
となる．両辺は $\pm1$ だから整数として等しい．
$p$ の法4の剰余によって指数の偶奇が決まり，
表示式を得る．
\end{proof}

\begin{theorem}[第二補充法則]\label{thm:second}
任意の奇素数 $p$ について
\[
 \left(\frac2p\right)=(-1)^{(p^2-1)/8}
 =
 \begin{cases}
 1&(p\equiv\pm1\pmod8),\\
 -1&(p\equiv\pm3\pmod8).
 \end{cases}
\]
\end{theorem}
\begin{proof}
$m=(p-1)/2$ とする．$1\le j\le m$ について
$2j$ は $2,\dots,p-1$ の偶数を動くので，
法 $p$ の最小正剰余は $2j$ 自身である．
$2j>p/2$ となるのは $j>p/4$ のときだから，
Gaussの補題の負符号の個数は
\[
 n(2)=m-\lfloor p/4\rfloor
\]
である．$p=8k+1,8k+3,8k+5,8k+7$ の
四つの場合に分けて計算すると，$n(2)$ の偶奇は
順に偶，奇，奇，偶となる．
一方 $(p^2-1)/8$ は整数であり，
その偶奇も順に偶，奇，奇，偶である．
よってGaussの補題から表示式を得る．
\end{proof}

\section{平方剰余の相互法則}

\begin{theorem}[平方剰余の相互法則]\label{thm:qr}
$p,q$ を異なる奇素数とする．
このとき
\[
 \left(\frac pq\right)\left(\frac qp\right)
 =(-1)^{\frac{p-1}{2}\frac{q-1}{2}}.
\]
すなわち，$p\equiv q\equiv3\pmod4$ の場合に限って
二つのLegendre記号の符号が異なる．
\end{theorem}
\begin{proof}
$m=(p-1)/2$，$n=(q-1)/2$ と置く．
Eisensteinの補題より
\[
 \left(\frac qp\right)
 =(-1)^{A},\qquad
 A=\sum_{i=1}^m\left\lfloor\frac{qi}{p}\right\rfloor,
\]
および
\[
 \left(\frac pq\right)
 =(-1)^{B},\qquad
 B=\sum_{j=1}^n\left\lfloor\frac{pj}{q}\right\rfloor
\]
である．
整数格子点
\[
 R=\{(i,j)\in\mathbb Z^2:
 1\le i\le m,\ 1\le j\le n\}
\]
を考える．$|R|=mn$ である．
各 $i$ に対して $qi/p<q/2$ なので
$\lfloor qi/p\rfloor\le n$ であり，
$A$ は $qi>pj$ を満たす格子点の個数に等しい．
同様に $B$ は $pj>qi$ を満たす格子点の個数に等しい．
$qi=pj$ が成立すると，
互いに異なる素数 $p,q$ は互いに素なので
$p\mid i$，$q\mid j$ となるが
$1\le i<p$，$1\le j<q$ に反する．
従って $R$ の各格子点は二種類のいずれか一方に属し，
$A+B=mn$ である．
よって
\[
 \left(\frac qp\right)\left(\frac pq\right)
 =(-1)^{A+B}=(-1)^{mn}
\]
となる．
\end{proof}

\begin{remark}
証明においては有限個の格子点を直線
$q x=p y$ の上下に分割した．
この分割の数え上げとGaussの補題の符号が，
平方剰余の相互法則を結びつける．
\end{remark}

\section{Jacobi記号}

\begin{definition}
$n>1$ を奇整数とし，
$n=\prod_{j=1}^rp_j^{e_j}$ を素因数分解する．
Jacobi記号を
\[
 \left(\frac an\right)
 =\prod_{j=1}^r\left(\frac a{p_j}\right)^{e_j}
\]
と定義する．$n=1$ のときは $(a/1)=1$ と約束する．
\end{definition}

\begin{proposition}\label{prop:jacobi}
$n>1$ を奇整数，$a,b\in\mathbb Z$ とすると，
\[
 \left(\frac{ab}{n}\right)
 =\left(\frac an\right)\left(\frac bn\right)
\]
が成り立つ．
また $\gcd(a,n)=1$ の場合には
$\left(\frac an\right)\in\{1,-1\}$ である．
\end{proposition}
\begin{proof}
Legendre記号の乗法性を各素因数 $p_j$ に適用し，
それぞれ $e_j$ 乗して掛け合わせれば乗法性を得る．
$\gcd(a,n)=1$ なら全ての $p_j$ で $p_j\nmid a$ だから，
各Legendre記号は $\pm1$ となり，
積も $\pm1$ である．
\end{proof}

\begin{theorem}[Jacobi記号の相互法則]\label{thm:jacobi-recip}
$m,n>0$ を互いに素な奇整数とする．このとき
\[
 \left(\frac mn\right)\left(\frac nm\right)
 =(-1)^{\frac{m-1}{2}\frac{n-1}{2}}.
\]
また任意の奇整数 $n>0$ に対し
\[
 \left(\frac{-1}{n}\right)=(-1)^{(n-1)/2},
 \qquad
 \left(\frac2n\right)=(-1)^{(n^2-1)/8}
\]
が成り立つ．
\end{theorem}
\begin{proof}
$m=\prod_ip_i^{a_i}$，$n=\prod_jq_j^{b_j}$ と分解する．
$m,n$ は互いに素なので $p_i\ne q_j$ である．
Legendre記号の相互法則を全ての $(p_i,q_j)$ に適用して
\[
 \left(\frac mn\right)\left(\frac nm\right)
 =(-1)^{
  \left(\sum_i a_i\frac{p_i-1}{2}\right)
  \left(\sum_j b_j\frac{q_j-1}{2}\right)
 }.
\]
奇整数 $u,v$ について
$\frac{uv-1}{2}\equiv\frac{u-1}{2}+\frac{v-1}{2}\pmod2$
が成り立つから，
右辺の指数は
$\frac{m-1}{2}\frac{n-1}{2}$ と法2で合同である．
従って相互法則を得る．
補充法則も素因子ごとのLegendre記号の公式を掛け合わせ，
奇整数の法4・法8の積の規則を用いれば従う．
\end{proof}

\begin{example}[Jacobi記号の値と平方剰余]
\[
 \left(\frac2{15}\right)
 =\left(\frac23\right)\left(\frac25\right)
 =(-1)(-1)=1
\]
だが，合同式 $x^2\equiv2\pmod{15}$ は整数解を持たない．
\end{example}
\begin{proof}
もし $x^2\equiv2\pmod{15}$ なら
特に $x^2\equiv2\pmod3$ となる．
しかし法3の平方剰余は $0,1$ のみなので矛盾する．
従ってJacobi記号が1であっても
平方剰余であるとは限らない．
\end{proof}

\section{相互法則の計算への応用}

\begin{example}[二次合同式の判定]
$p$ を $5$ と異なる奇素数とする．
合同式 $x^2\equiv5\pmod p$ が解を持つための
必要十分条件は
$p\equiv1,4\pmod5$ である．
\end{example}
\begin{proof}
$5\equiv1\pmod4$ なので平方剰余の相互法則から
$\left(\frac5p\right)=\left(\frac p5\right)$ である．
法5の非零平方剰余は $1,4$ だから，
右辺が1になることと $p\equiv1,4\pmod5$ は同値である．
Legendre記号の定義によって結論が従う．
\end{proof}

\begin{example}[$-1$ と $2$ の合同式]
奇素数 $p$ について，
$x^2\equiv-1\pmod p$ は
$p\equiv1\pmod4$ の場合に限って解を持つ．
また $x^2\equiv2\pmod p$ は
$p\equiv1,7\pmod8$ の場合に限って解を持つ．
\end{example}
\begin{proof}
前者は第一補充法則，後者は第二補充法則による．
\end{proof}

\begin{remark}
平方剰余の相互法則は，法の異なる二次合同式を結びつけ，
素数を法とする平方剰余判定を相互に変換する．
その証明は有限群・合同式・格子点の数え上げを横断している．
\end{remark}


\section{到達課題の完全解答}
\noindent\textbf{平方剰余相互法則を使い切る．}
\begin{proposition}
$3$ は $13$ と $11$ を法として平方剰余であるが，
$7$ を法として平方非剰余である．
\end{proposition}
\begin{proof}
$4^2=16\equiv3\pmod{13}$，
$5^2=25\equiv3\pmod{11}$ だから前二者は平方剰余である．
$7$ を法とする非零平方剰余は
$1^2,2^2,3^2$ の剰余 $1,4,2$ のみであり，
$3$ は含まれない．従って非剰余である．
\end{proof}
\begin{proposition}[相互法則による別の判定]
$\bigl(\frac5{11}\bigr)=1$ である．
\end{proposition}
\begin{proof}
$5\equiv1\pmod4$ なので相互法則より
$\left(\frac5{11}\right)
=\left(\frac{11}5\right)=\left(\frac15\right)=1$．
\end{proof}
\noindent\textbf{到達確認．}
$\left(\frac3{7}\right)$ を相互法則から求めよ．
\par\noindent\textbf{解答．}
$3\equiv7\equiv3\pmod4$ だから
$\left(\frac3{7}\right)
=-\left(\frac73\right)
=-\left(\frac13\right)=-1$ である．
\end{document}
